\section{Licencias Software}

Las tecnologías y herramientas utilizadas para el desarrollo de este proyecto se encuentran bajo las siguientes licencias.

\subsection{Licencia MIT - Massachusetts Institute of Technology}

La licencia MIT~\cite{mit_license} es una licencia \textit{software} procedente del Instituto de Tecnología de Massachusetts. Esta cede sin restricciones los derechos de: uso, copia, modificación y distribución. Además, existe una única condición, incluir la nota de \textit{copyright} y los derechos anteriormente mencionados en cualquier copia del código distribuido, pudiendo verse invalidada la licencia en caso de no cumplimiento.

\subsection{Licencia GPLv2 - \textit{General Public License version 2}}

También conocida como \textit{Licencia Pública General - GNU}~\cite{gplv2_license} es una licencia de derecho de autor utilizada en \textit{software} libre y código abierto. Su principal objetivo es el de permitir la distribución de \textit{software} libre y protegerlo legalemente de cualquier intento de apropiación que impida su uso a nuevos usuarios. La licencia permite los derechos de: uso, estudio, modificación y distribución. Por otro lado, entre las condiciones se encuentran: disponibilidad del código fuente, conservación de avisos legales, distribución del software sin garantía y \textit{copyleft}.

\subsection{Licencia LPPL - \textit{LaTex Project Public License}}
La LPPL, tambien conocida como \textit{Licencia Publica del Proyecto LaTex}~\cite{latex_license}, es una licencia de código gratuito. Entre los derechos que la licencia otorga, encontramos: uso, distribución y modificación. Por otro lado, respecto a las condiciones se distinguen: identificación clara de versiones modificadas, conservación de avisos legales, copia de la licencia en cualquier distribución, distribución del software sin garantía y \textit{copyleft}.

\subsection{Licencia Apache}
La licencia Apache~\cite{apache_license}, es una licencia de software libre creada por la \textit{Apache Software Foundation} - ASF. Proporciona derechos de distribución y modificación, con una única condición, la existencia de dos archivos, una copia de la licencia y un aviso.

\subsection{Licencia EULA - \textit{End User License Agreement}}
Se trata de un acuerdo de licencia de usuario final~\cite{eula_license}, un contrato donde actuan un proveedor y un cliente. Este acuerdo de licencia especifica los derechos y restricciones sobre el software que afectan al usuario.

\subsection{Licencias Comerciales}
Algunos programas utilizados llevan atribuida una licencia comercial. Esta concede el derecho a usar los programas siempre y cuando se cumplan los términos del servicio respectivo.